% Chapter 2, Figure 15: a step disturbance of intensity M_s (scan: figures/ch2/fig15.png).
% Curve: fig15.py, the step input as the text defines it (unnumbered display before eq. (27)): f_x(t) = 0 for
% t < 0 (along the t axis), M_s for t >= 0, the riser on the vertical axis at t = 0, as in the scan.
% Four-quadrant axes as printed: the t axis runs as far left of the origin as right, the vertical axis well below 0.
% Same axis size as Figs 16-19.
\documentclass[11pt]{standalone}
\usepackage{tamrfig}
\begin{document}
\begin{tikzpicture}
  \begin{axis}[tamr sketch,
      width=3.8in, height=2.5in,
      xmin=-1, xmax=1.08, ymin=-1.24, ymax=1.3,
      ytick={1}, yticklabels={$M_s$},
      xlabel={$t$ (sec)}, ylabel={Yawing moment $M_X$ (dyn-cm)},
      % tamr sketch's rotate=-90 turns this label to read downward: reset it to upright, above the arrow
      every axis y label/.style={at={(ticklabel* cs:1.0)}, anchor=south}]
    \addplot[series1] table[x=t, y=M, col sep=comma] {figures/v2/ch2/fig15.csv};
    \node[anchor=north west, font=\footnotesize, text=ink2] at (axis cs:0,0) {0};
  \end{axis}
\end{tikzpicture}
\end{document}
